% ---- pg-590 （书页 578） ----

因此, 将高风险申请人划分为“低”(接受)的代价是将低风险申请人划分为“高”(拒绝)的代价的 5 倍. 在实践中, 接受一个低的信用风险应该产生收益, 或负成本. 为满足此条件, 我们从成本矩阵的第一行减去 1 美元, 以获得更符合“实际”的成本矩阵如下:

\begin{mvtable}
\setlength{\tabcolsep}{14pt}%
\begin{tabular}{cc|cc}
\toprule
\multicolumn{2}{c|}{} & \multicolumn{2}{c}{预测(决策)} \\
\cmidrule(lr){3-4}
\multicolumn{2}{c|}{} & 低(接受) & 高(拒绝) \\
\midrule
\multirow{2}{*}{实际} & 低 & $-1$ 美元 & 0 \\
& 高 & 5 美元 & 0 \\
\bottomrule
\end{tabular}
\end{mvtable}

该成本矩阵与原来的矩阵产生相同的决策, 但是其结果从期望成本目标函数的角度更容易解释. 例如, 经过进一步的调整, 负的期望成本得分可能意味着潜在的收益, 因此申请人可能是低信用风险.

然后, 应该加工(变换)输入变量, 指定模型(或算法), 对流程图中每个图标设置需要的参数. 此时通过点击流程图中的任一图标可以执行程序到该点. 例如, 点击 Score 图标可以使程序从头执行到该图标(含). 在每个被执行的图标中的结果可以通过点击该图标得到.

我们用提升图解释模型评估. 这些提升图可从 Assessment 图标中得到.

考虑逻辑斯蒂回归分类. 逻辑斯蒂回归函数由训练数据确定, 用此函数对验证集的每个案例计算期望成本. 这些期望成本“得分”可以从小到大重新排序, 并且按照 10,20,$\cdots$ 百分位数分组. 第一个百分位组有 42(419 的 10\%)个申请人, 他们有最小的负期望成本(最大的潜在收益); 第二个百分位组同样有 42(419 的 10\%)个申请人, 如此类推. (从分类的观点看, 那些有负期望成本的申请人应被分入低风险, 而有正期望成本的申请人应被分入高风险.)

如果模型没有预测能力, 我们可以期望低信用风险在各个组内大致是均匀分布的. 也就是说, 在每个百分位组内的 42 个人中, 大致有 10\% 或 $0.10\times299=30$ 个低信用风险.

当验证数据被评分后, 我们可以对每个百分位组内的低风险计数. 例如, 在第一个百分位组的 42 个申请人中, 有 40 个是低风险, 其“捕获的响应率”(captured response rate)为 $40/299=0.133$ 或 13.3\%. 在这种情况下, 第一百分位组的提升为预测的低风险数目与随机指派的低风险数目的比, 即
\[
\mathrm{Lift}=\frac{40}{30}=1.33
\]
提升值说明模型比随机指派多指派 $10/299=0.33$ 或 3.3\% 的低风险到第一百分位组(最大负期望成本组).\textsuperscript{①}

提升统计量可表示为个别值(非累计的)或累计值. 例如, 在逻辑斯蒂回归分类中, 第二百分位组也有 40 个低风险. 因此前两个百分位组累计提升为
\[
\mathrm{Lift}=\frac{40+40}{30+30}=1.33
\]
\hyperref[fig:12.31]{图 12.31} 是逻辑斯蒂回归的累计提升图.

对分类树工具和神经网络工具也可得到提升与累计提升统计量. 对每种分类方法, 对整个数据集记分(计算期望成本), 申请人按照得分从小到大排序, 并创建百分位组. 提升计算按照对逻辑斯蒂回归的计算进行. 所有三个累计提升图表示在\hyperref[fig:12.32]{图 12.32} 中.

\vspace{6pt}
\noindent\rule{\linewidth}{0.4pt}
\par\vspace{2pt}
{\zihao{6}\noindent ①\quad 由于有舍入误差, 这里计算的提升值与后面提升图中的数值有些微小差异.\par}
